\documentclass[10pt,a4paper]{amsart}
\usepackage[margin=19mm]{geometry}
\usepackage{amsmath,amssymb,mathtools}
\usepackage[scr=rsfs,cal=euler]{mathalfa}
\usepackage{xfrac}
\usepackage[unicode,hidelinks,bookmarks=false]{hyperref}
\newcommand{\ie}{i.e.\ }
\newcommand{\cf}{cf.\ }
\newcommand{\resp}{respectively\ }
\newcommand{\Subsection}{\subsection}
\providecommand{\nobreakdash}{\nobreak\hbox{-}}
\setlength{\emergencystretch}{2em}
\multlinegap0pt
\usepackage{xcolor}
\usepackage{amssymb}
\usepackage{tcolorbox}
\newtheorem{theorem}{Theorem}
\newcommand\Ren[1]{\mathbb{R}^{#1}}
\newcommand\ar{\mathbb{R}}
\newcommand{\inverseCT}{{\mathcal T}^{-1}}
\newcommand\inverseRadonPart[2]{\left(\frac{1}{{#1}}\frac{d}{d{#1}}\right)^{#2}}
\newcommand\smallcirc{\scriptstyle\circ}
\title{On rotationally-symmetric norms}
\author{Yossi Lonke}
\date{Source: arXiv:2609.14554v1 (CC BY 4.0)}
\begin{document}
\maketitle

\noindent\emph{Source and licence.} On rotationally-symmetric norms. Version arXiv:2609.14554v1 (CC BY 4.0), distributed by arXiv under the Creative Commons Attribution 4.0 International licence, http://creativecommons.org/licenses/by/4.0/, as recorded in the arXiv licence metadata for this exact version and shown on the abstract page https://arxiv.org/abs/2609.14554v1. Full licence notice: \texttt{licenses/arxiv-2609.14554-CC-BY-4.0.txt}.\par\smallskip

\noindent\textit{Editorial scope.} Counted unit: the article's theorem myInversionTheorem (source0.tex lines 703--711). The article's complete original proof of it is delivered with it (source0.tex lines 712--781, one \texttt{proof} environment of 70 lines). No mathematics is written, repaired or re-derived: the statement and the proof are the article's own lines, in the article's own notation, and no retained line is substituted or re-flowed.  Retained uncounted as the source-local context the proof needs: lines 693--695.  Nothing was omitted from any retained span, so the package carries no omission record.  No retained line contains a citation, so the wrapper carries no bibliography and no bibliography entry.  No retained line contains a figure environment, an image-inclusion command or drawing code, so no image is delivered and the package declares no asset dependency.  Editorial formatting only, in addition: an AMS \texttt{amsart} preamble that restates the article's own declarations and macros used by the retained lines, namely the environments \texttt{theorem} on separate counters, together with the standard packages those lines need.  The retained environments print in this document as Theorem 1 in order of appearance, whereas the article prints the same items with the numbering of its own full document.  The complete native source of the article as distributed in the arXiv e-print for this exact version is bundled unchanged as \texttt{sources/210332/source0.tex}.
\begin{equation}\label{differentialOperator}
c_n\left [(1-x^2)\frac{d^2}{dx^2}-(n-1)x\frac{d}{dx}+n-1\right]\circ x\left (\frac{1}{x}\frac{d}{dx}\right)^{\frac{n}{2}-1}.
\end{equation}

\begin{theorem}\label{myInversionTheorem}
Let $f\in C^{\infty}(\Ren{2n}\backslash\{0\})$ be a rotationally symmetric, continuous, and homogeneous of degree~$1$ real valued function. Put
\[ \varphi(t)=f(e_{2n-1}+te_{2n}),\qquad (t\in\ar). \]
Then the rotationally symmetric inverse cosine transform  $\inverseCT f$ restrticed to  the inverval $(0,1)$ is given by the formula 

\begin{equation}\label{myInversionFormula}
 \alpha_{2n}^{-1}(\inverseCT f)=\frac{(-2)^{n-1}}{x^{2n+1}}\frac{d^{n-1}}{dt^{n-1}}(\varphi''(\sqrt{t}))\Big|_{t\to (\frac{1-x^2}{x^2})}.
\end{equation}
\end{theorem}

\begin{proof} 
\begin{comment}We shall prove the formula for every $\varphi\in C^{\infty}(0,\infty)$ and deduce by duality that it also
holds for the dual space of distributions.
\end{comment}

For every $k$-times differentiable function $f$ defined in $\ar^+$,
\begin{equation}\label{keyObservation}
\frac{1}{2^k}\left(\frac{1}{x}\frac{d}{dx}\right)^k(f)=\frac{d^k}{dt^k}(f{\smallcirc}\sqrt{t})\Big|_{t\to x^2},
\end{equation}
as can easily be verified by induction on $k$.

\medskip

Since $f$ is homogeneous of degree $1$,  we have 
\begin{equation}\label{homogeneity}
 x^{2n-3}f(\sin^{-1}x)=x^{2n-2}\varphi(\frac{\sqrt{1-x^2}}{x}),\qquad (0<x <1).
 \end{equation}
For $0<t<1$, put  $h(t)=(\frac{1-t}{t})^{1/2}$. In view of (\ref{keyObservation}) and (\ref{homogeneity}), we can write
\begin{equation}\label{twoDiffops}
\frac{1}{2^{n-1}}\inverseRadonPart{x}{n-1}(x^{2n-3}f(\sin^{-1} x)) = \frac{d^{n-1}}{dt^{n-1}}(t^{n-1}\varphi{\smallcirc} h(t))\Big|_{t\to x^2}.
\end{equation}
For $m\in\mathbb{N}$, define a differential operator on $C^{\infty}(0,\infty)$ by
\[ F_m(g)=(1-x^2)\frac{d^2g}{dx^2}-(m-1)x\frac{dg}{dx}+(m-1)g,\]
and for every $m\in\mathbb{N}$ and $g\in C^{\infty}(0,\infty)$, define a function of $0<t<1$ by
\[ H_m^g(t)=\frac{d^{n-1}}{dt^{n-1}}(t^{n-1}g{\smallcirc} h). \]
By (\ref{differentialOperator}) and (\ref{twoDiffops}),  for $0<x<1$ the formula (\ref{myInversionFormula})  is 
equivalent to 

\begin{equation}\label{eq1}F_{2n}(x H_n^{\varphi}(x^2)))=
\frac{(-1)^{n-1}}{x^{2n+1}}
\frac{d^{n-1}}{dt^{n-1}}
(\varphi''(\sqrt{t}))\Big|_{t\to h^2(x^2)}
\end{equation}

\medskip

An application of the differential operator $F_{2n}$ to the expression $xH_n^{\varphi}(x^2)$ is straightforward:
\begin{equation}\label{eq2}
F_{2n}(xH_n^\varphi(x^2))=(H_n^\varphi)'(x^2)(6x-4x^3(n+1))+4x^3(1-x^2)(H_n^\varphi)''(x^2).
\end{equation}

It remains to show that the right hand sides of (\ref{eq1}) and (\ref{eq2}) are equal. If~${n=1}$, a routine computation verifies the equality.

\medskip


\begin{comment}For $n\geq 2$, we first prove that the right hand sides of (\ref{eq1}) and (\ref{eq2})  are equal under the assumption that $\varphi(x)=x^m$ for some $m\in\mathbb{N}$. By linearity, we deduce that they are equal whenever $\varphi$ is a polynomial.  Then, we restrict the values of $x$ to a compact interval $[a,b]\subset (0,1)$. The smooth bijection $h(x)$ maps it onto $[0,1/a]$ and the polynomials
are dense in $C^{\infty}[0,1/a]$ in the topology of uniform convergence of all derivatives.\textcolor{orange}{double check that}  By continuity of the differential
operators in this topology, we finally deduce that the formula holds for all $\varphi\in C^{\infty}(0,\infty)$ and $0<x<1$, hence it holds also in the distribution sense.
\end{comment}

For every $0<a<b<1$, the right hand sides of  (\ref{eq1}) and (\ref{eq2}) are  both linear in~$\varphi$ and continuous in 
$C^{\infty}[a,b]$. Since the set of polynomials is dense in $C^{\infty}[a,b]$ with the topology of uniform convergence of all derivatives, it suffices to check that (\ref{eq1}) and (\ref{eq2})  agree on all monomials $\varphi(x)=x^m$, where $m\in\mathbb{N}$. Consider two cases:
\begin{enumerate}
\item[1.] $m=2l$ is an even integer and $l<n$. Hence $\varphi''(\sqrt{x})=2l(2l-1)x^{l-1}$. Since $l-1<n-1$, the right hand side of (\ref{eq1}) vanishes. On the other hand,
\[ x^{n-1}\varphi(h(x))=x^{n-l-1}(1-x)^l \]
is a polynomial of degree $n-1$, so with $\varphi=x^{2l}$, the function $H_n^\varphi$ is  constant, and the right hand side of (\ref{eq2}) vanishes as well.
\item[2.] In the second case,  $m$ is either an even integer at least $2n$, or else an odd natural number. With $\varphi=x^m$, expand the right hand side of 
\[ H_n^\varphi(x)=\frac{d^{n-1}}{dx^{n-1}}(x^{n-1-m/2}(1-x)^{m/2}), \]
to obtain
\[ (-1)^{n-1}\sum_{k=0}^{n-1}{n-1\choose k}(-1)^k(\frac{1-x}{x})^{m/2-n+k+1}(n-1-m/2)^{(k)}(m/2)^{(n-k+1)},\]
where  $(a)^q=a(a-1)\dots (a-q+1)$ for real $a$ and $q\in\mathbb{N}$. 
A substitution of this sum for $H_n^\varphi(x)$ into the right hand side of (\ref{eq2}) results in a complicated but elementary calculation that yields the expression 
\begin{equation}\label{final}
\frac{(-1)^{n-1}}{x^{m+1}}m(m-1)(1-x^2)^{m/2-n}(m/2-1)^{(n-1)}, 
\end{equation}
which for  $\varphi(x)=x^m$, is precisely the right hand side of (\ref{eq1}). 
\end{enumerate}
This holds for every compact interval $[a,b]\subset (0,1)$ and therfore holds for every ${\varphi\in C^{\infty}(0,1)}$.
\end{proof}

\end{document}
